% TEMPLATE: step chart (a quantity that narrows through stages, held
% constant between them and jumping at each boundary -- never a smooth
% interpolation between categorical stages). The stage that narrows the
% most is the chapter's claim, marked with a focus datum and a direct label
% for its drop; every level is also on the tick list, so no legend is
% needed to read a count off the axis.
\begin{figure}[H]
\centering
\pdfigurehue{pdcobalt}%
\begin{tikzpicture}[pd figure]
\begin{axis}[pd axis,
  width=8.0cm,height=4.3cm,
  xmin=-0.3,xmax=3.5,ymin=0,ymax=520,
  xtick={0,1,2,3},
  xticklabels={submitted,triaged,reviewed,admitted},
  ytick={96,150,340,480},
  ylabel={candidates remaining}]
  \addplot[pd series,const plot,mark=none] coordinates {(0,480)(1,340)(2,150)(3,96)(3.4,96)};
  \node[pd focus datum] at (axis cs:2,150) {};
  \node[pd label] at (axis cs:1.32,410) {$-140$};
  \node[pd focus label] at (axis cs:2.32,245) {$-190$};
  \node[pd label] at (axis cs:3.18,140) {$-54$};
\end{axis}
\end{tikzpicture}
\caption{A step chart of candidates remaining at each review stage: the
count holds within a stage and jumps at the boundary, never interpolates
between two categorical stages. The review stage removes the most
candidates, $-190$, roughly a third more than triage's $-140$ and four
times admission's $-54$; every stage's count is also on the tick list, so
each level can be read directly off the axis. [internal, illustrative]}
\label{fig:tpl-step-chart}
\end{figure}
